% Time Integral

\subsubsection{Infinite-range screened Coulomb potentials}
%%%%%%%%%%%%%%%%%%%%%%%%%%%%%%%%%%%%%%%%%%%%%%%%%%%%%%%%%%
%
The integral in \cref{eq:sn_reduced_tau} for the calculation of the time integral is also of the form \cref{eq:sn_integration_integral_R}. Applying the same substitution of variables \cref{eq:sn_transform1} to \cref{eq:sn_reduced_tau}, one obtains
%
\begin{equation}
    \label{eq:sn_tau_u}
    T = R_0 - 2 R_0 \int_0^1
        \left[ \frac{ 1 }{ g \left( \frac{ R_0 }{ 1-u^2 } \right) } - 1 \right]
        \frac{ u }{ (1-u^2)^2 }
        \, \mathrm{d}u
\end{equation}
%
with $g(\frac{R_0}{1-u^2})$ as given by \cref{eq:sn_cmangl_legendre_g_u}.
Combining the two terms in brackets in \cref{eq:sn_tau_u} and multiplying numerator and denominator with $1 + g(\frac{R_0}{1-u^2})$ gives
%
\begin{equation}
    \label{eq:sn_tau_u1}
    T = R_0 - 2 R_0 \int_0^1
        \frac{ 1 - g^2 \left( \frac{ R_0 }{ 1-u^2 } \right) }{
               g \left( \frac{ R_0 }{ 1-u^2 } \right)
               \left[ 1 + g \left( \frac{ R_0 }{ 1-u^2 } \right) \right] }
        \frac{ u }{ (1-u^2)^2 }
        \, \mathrm{d}u .
\end{equation}
%
Inserting \cref{eq:sn_cmangl_legendre_g_u} in the numerator and \cref{eq:sn_cmangl_legendre_g_G} in the demoninator, and using \crefrange{eq:sn_cmangl_legendre_G_chi}{eq:sn_cmangl_legendre_rho} yields
%
\begin{equation}
    \label{eq:sn_tau_G}
    T = R_0 - 2 P \int_0^1
        \frac{ 1 + \rho \, \frac{ 1 }{ 1-u^2 } \, \Phi \left( \frac{R_0}{1-u^2} \right) }{
               G(u) \left[ 1 + u \frac{P}{R_0} G(u) \right] }
        \, \mathrm{d}u ~.
\end{equation}
%
Note that the term $\frac{ 1 }{ 1-u^2 } \, \Phi \left( \frac{R_0}{1-u^2} \right)$ does not diverge if $R\, \Phi(R) \rightarrow 0$ as $R \rightarrow \infty$, in other words, if $\Phi(R)$ decays faster that $1/R$.

Using Gauss-Legendre quadrature, the time integral may be approximated by
%
\begin{equation}
    \label{eq:sn_tau_G_legendre}
    T = R_0 - 2 P \sum_{i-1}^n \tilde{w}_i
        \frac{ 1 + \rho \, \frac{ 1 }{ 1-\tilde{u}_i^2 } \, \Phi \left( \frac{R_0}{1-\tilde{u}_i^2} \right) }{
        G(\tilde{u}_i) \left[ 1 + \tilde{u}_i \frac{P}{R_0} G(\tilde{u}_i) \right] } ~.
\end{equation}
%
with $\tilde{u}_i$ and $\tilde{w}_i$ defined in \cref{eq:sn_gauss_legendre_transform_infinite}.
In terms of the function $H(u)$, \cref{eq:sn_H_def}, \cref{eq:sn_tau_G_legendre} reads
%
\begin{equation}
    \label{eq:sn_tau_H_legendre}
        T = R_0 - 2 \sum_{i-1}^n \tilde{w}_i
        \frac{ P^2 + \frac{R_0}{\varepsilon} \, \frac{ 1 }{ 1-\tilde{u}_i^2 } \, \Phi \left( \frac{R_0}{1-\tilde{u}_i^2} \right) }{
        H(\tilde{u}_i) \left[ 1 + \frac{\tilde{u}_i}{R_0} H(\tilde{u}_i) \right] } ~.
\end{equation}
%
which explicitly shows that the second term does not vanish for $P \rightarrow 0$.

\subsubsection{Finite-range screened Coulomb potentials}
%%%%%%%%%%%%%%%%%%%%%%%%%%%%%%%%%%%%%%%%%%%%%%%%%%%%%%%%
%
For finite-range screened Coulomb potentials we are interested only in $P < R_\mathrm{max}$ since for $\theta = 0$ the outgoing asymptote is identical with the incoming asymptote independent of the time integral.
Like for the scattering angle, we split the integral in \cref{eq:sn_reduced_tau} in two parts:
%
\begin{equation}
    \label{eq:sn_reduced_tau_finite}
    T = R_0 -
        \int_{R_0}^{R_\mathrm{max}} \left[ \frac{ 1 }{ g(R) } - 1 \right] \mathrm{d}R -
        \int_{R_\mathrm{max}}^\infty \left[ \frac{ 1 }{ g(R) } - 1 \right] \mathrm{d}R ~.
\end{equation}
%
The second integral may be evaluated analytically using \cref{eq:sn_g_phi=0}
%
\begin{equation}
    T = R_0 -
        (R_\mathrm{max} - \sqrt{ R_\mathrm{max}^2 - P^2 }) -
        \int_{R_\mathrm{max}}^\infty \left[ \frac{ 1 }{ g(R) } - 1 \right] \mathrm{d}R ~.
\end{equation}

Applying the substitution \cref{eq:sn_transform1} to the remaining integral, using the identity
%
\begin{equation*}
    \frac{1}{g} - 1 = \frac{1-g}{g} = \frac{1-g^2}{g(1+g)}
\end{equation*}
%
and Eqs.~(\ref{eq:sn_cmangl_legendre_g_u}), (\ref{eq:sn_cmangl_legendre_g_G}), (\ref{eq:sn_cmangl_legendre_G_chi}), and (\ref{eq:sn_cmangl_legendre_rho})
we obtain
%
\begin{equation}
    \label{eq:sn_reduced_tau_finite_G_legendre}
    T = R_0 -
        (R_\mathrm{max} - \sqrt{ R_\mathrm{max}^2 - P^2 }) - 2 P \int_0^{u_\mathrm{max}}
        \frac{ 1 + \rho \, \frac{ 1 }{ 1-u^2 } \, \Phi \left( \frac{R_0}{1-u^2} \right) }{
            G(u) \left[ 1 + u \frac{P}{R_0} G(u) \right] }
        \, \mathrm{d}u \\
\end{equation}
%
with $u_\mathrm{max}$ as defined in \cref{eq:sn_umax}.

Using Gauss-Legendre quadrature, the time integral may finally be approximated by
%
\begin{equation}
    \label{eq:sn_cmangl_legendre_tau_chi_finite}
    T = R_0 - R_\mathrm{max} + \sqrt{ R_\mathrm{max}^2 - P^2 } -
    2 P \sum_{i=1}^n \tilde{w}_i
    \frac{ 1 + \rho \, \frac{ 1 }{ 1-\tilde{u}_i^2 } \, \Phi \left( \frac{R_0}{1-\tilde{u}_i^2} \right) }{
        G(\tilde{u}_i) \left[ 1 + \tilde{u}_i \frac{P}{R_0} G(\tilde{u}_i) \right] } ~.
\end{equation}
%
with $\tilde{u}_i$ and $\tilde{w}_i$ defined in \cref{eq:sn_gauss_legendre_transform_finite}.
In terms of the function $H(u)$, \cref{eq:sn_H_def}, \cref{eq:sn_cmangl_legendre_tau_chi_finite} reads
%
\begin{equation}
    \label{eq:sn_cmangl_legendre_tau_H_finite}
    T = R_0 - R_\mathrm{max} + \sqrt{ R_\mathrm{max}^2 - P^2 } -
        2 \sum_{i=1}^n \tilde{w}_i
        \frac{ P^2 + \frac{R_0}{\varepsilon} \, \frac{ 1 }{ 1-\tilde{u}_i^2 } \, \Phi \left( \frac{R_0}{1-\tilde{u}_i^2} \right) }{
        H(\tilde{u}_i) \left[ 1 + \frac{\tilde{u}_i}{R_0} H(\tilde{u}_i) \right] } ~.
\end{equation}


\subsubsection{Convergence}
%%%%%%%%%%%%%%%%%%%%%%%%%%%
%
\begin{figure}[htbp]
    \centering
    \includegraphics[scale=0.5]{error_legendre_tau_eps1e-4_zbl.pdf} ~~
    \includegraphics[scale=0.5]{error_legendre_tau_eps1e-4_nlhlin.pdf} \\ \smallskip
    \includegraphics[scale=0.5]{error_legendre_tau_eps0.1_zbl.pdf} ~~
    \includegraphics[scale=0.5]{error_legendre_tau_eps0.1_nlhlin.pdf} \\ \smallskip
    \includegraphics[scale=0.5]{error_legendre_tau_eps100_zbl.pdf} ~~
    \includegraphics[scale=0.5]{error_legendre_tau_eps100_nlhlin.pdf}
    \caption{Relative errors in the time integral \(T = \tau/a_\mathrm{I}\) of n-point Gauss-Legendre quadrature as a function of reduced impact parameter for three energies using \cref{eq:sn_cmangl_legendre_tau_chi_finite} with the number of abscissae $n=1$,2, 3, 4, 5, and 10.
        As interatomic potentials the universal ZBL potential (left) and the NLHlin potential for $Z_1=33$ and $Z_2=14$ were used together with a screening length $a_\mathrm{I}$ given by \cref{eq:sn_aNLH}.
        The reference simulations were evaluated using \cref{eq:sn_cmangl_legendre_tau_chi_finite} with 32-point Gauss-Legendre quadrature.
        The relative accuracy of the apsis is better than \(10^{-3}\).}
    \label{fig:error-legendre-tau}
\end{figure}
%
In \cref{fig:error-legendre-tau}, the relative errors in \(T = \tau/a_\mathrm{I}\) as calculated by \cref{eq:sn_cmangl_legendre_tau_chi_finite} are shown as a function of reduced impact parameter $P$ for reduced energies \(\varepsilon = 10^{-4}\), \(\varepsilon = 10^{-1}\), and \(\varepsilon = 100\), and different numbers of abscissae $n$ in the Gauss-Legendre quadrature.
The results for the ZBL potential are shown in the left column, while those for the NLHlin potential are found in the right column.
An accuracy of at least 1\% is reached with $n = 4$.
For the NLHlin potential, convergence is even faster in general.

In summary, 4-point Gauss-Legendre quadrature allows the calculation of the scattering angle with an accuracy of  better than 10$^{-3}$ and the calculation of the time integral with an accuracy of better than 1\%.
Using the same number of abscissae for the calculation of both quantities has the advantage that the $\Phi \left( \frac{R_0}{1-u^2} \right)$ values evaluated for the scattering angle can be reused for the time integral.

